% --- 1. INFORMAZIONI GENERALI ---
\section{Informazioni Generali}
\begin{itemize}
    \item \textbf{Luogo/Piattaforma:} Server Discord\textsubscript{G}
    \item \textbf{Data:} 2026-06-30
    \item \textbf{Orario di inizio:} 14:30
    \item \textbf{Orario di fine:} 15:40
    \item \textbf{Responsabile\textsubscript{G}:} Armando Moda Scarati
    \item \textbf{Scriba:} Sofia De Blasi
\end{itemize}

\vspace{0.5cm}
\textbf{Partecipanti presenti:}
\begin{table}[ht]
    \centering
    \renewcommand{\arraystretch}{1.2}
    \begin{tabularx}{\textwidth}{X l}
        \toprule
        \textbf{Nome e Cognome} & \textbf{Matricola} \\
        \midrule
        Sofia De Blasi & 2111014 \\
        Roberto Mariano Doroftei & 2111031 \\
        Filippo Idiometri & 2035617 \\
        Armando Moda Scarati & 2082864 \\
        Anton Ricardo Rupi & 2054304 \\
        \bottomrule
    \end{tabularx}
\end{table}

\vspace{0.5cm}
\textbf{Partecipanti assenti:}
\begin{table}[ht]
    \centering
    \renewcommand{\arraystretch}{1.2}
    \begin{tabularx}{\textwidth}{X l}
        \toprule
        \textbf{Nome e Cognome} & \textbf{Matricola} \\
        \midrule
        Francesco Marcon & 2110990 \\
        \bottomrule
    \end{tabularx}
\end{table}

% --- 2. ORDINE DEL GIORNO ---
\section{Ordine del Giorno}

Gli argomenti previsti per la riunione odierna sono:

\begin{enumerate}
	\item Sprint\textsubscript{G} review e retrospective.
	\item Allineamento sulle automazioni CI/CD\textsubscript{G} e metriche di codice e test.
	\item Stesura documenti (Verbale Interno, Diario di Bordo\textsubscript{G}, PdP, PdQ).
	\item Applicazione delle correzioni sull'AdR ed e-mail per Zucchetti\textsubscript{G}.
	\item Pianificazione dello sprint e assegnazione ruoli.
\end{enumerate}


% --- 3. RESOCONTO E DISCUSSIONE ---
\section{Resoconto della Riunione}

\subsection{Sprint review\textsubscript{G} e retrospective}
Il gruppo ha effettuato la verifica dello sprint appena concluso, constatando che tutte le attività previste sono state ultimate correttamente. Durante l'analisi retrospettiva, è emerso un rallentamento nel lavoro di preparazione alla Product Baseline\textsubscript{G} (PB\textsubscript{G}), causato principalmente dall'intensificarsi degli impegni universitari e dalla sessione d'esami dei membri. Come azione correttiva, il team ha deciso di sfruttare il periodo di stallo in attesa dei colloqui e dei ricevimenti ufficiali per concentrare il tempo disponibile sulle attività più imminenti. Nello specifico, ci si focalizzerà sulle comunicazioni aziendali e sulle attività gestionali (stesura verbali, aggiornamento PdP e PdQ).

\subsection{Aggiornamenti operativi e comunicazioni aziendali}
Si è discusso della necessità di apportare correzioni all'Analisi dei Requisiti\textsubscript{G}, definendo un'attività di "palestra" per comprenderle appieno e preparare le relative domande da sottoporre all'azienda proponente. Questa attività formativa non impatterà sul monte ore rendicontato. È stato inoltre stabilito di redigere un'e-mail da inviare a Zucchetti per finalizzare una data per il colloquio. Il team ha poi proseguito con l'organizzazione delle attività per il setup dell'ambiente di Continuous Integration/Continuous Deployment (CI/CD) e la selezione delle metriche necessarie per la valutazione della qualità del codice e dei test. 

\subsection{Pianificazione dello sprint successivo e assegnazione dei ruoli}
Il gruppo ha definito l'assegnazione dei ruoli per garantire una copertura adeguata delle task\textsubscript{G} previste nel breve termine, mantenendo basso il carico orario ove necessario a causa degli esami:
\begin{itemize}
    \item \textbf{Responsabile:} Armando Moda Scarati, incaricato di redigere il verbale interno e il Diario di Bordo.
    \item \textbf{Analista\textsubscript{G}:} Sofia De Blasi, con l'obiettivo di svolgere attività non rendicontate (palestra) per la comprensione delle correzioni da applicare all'AdR e la stesura delle domande per Zucchetti.
    \item \textbf{Amministratori (divisi in 3 flussi):} Roberto Mariano Doroftei si occuperà dell'aggiornamento del PdQ e della stesura dell'e-mail per Zucchetti; Filippo Idiometri seguirà il setup delle automazioni e delle metriche; Francesco Marcon avrà il compito\textsubscript{G} di aggiornare il PdP.
    \item \textbf{Verificatore\textsubscript{G}:} Anton Ricardo Rupi.
\end{itemize}

% --- 4. DECISIONI PRESE ---
\section{Decisioni Prese}

\renewcommand{\arraystretch}{1.3}
\vspace{-0.3cm}
\noindent
\begin{longtable}{c p{12cm}}
    \toprule
    \textbf{Codice} & \textbf{Decisione} \\
    \midrule
    \endfirsthead

    \toprule
    \textbf{Codice} & \textbf{Decisione} \\
    \midrule
    \endhead

    \midrule
    \multicolumn{2}{r}{\textit{Continua nella pagina successiva}} \\
    \midrule
    \endfoot

    \bottomrule
    \endlastfoot

    \textbf{VI-20.1} & {Ricalibrazione del carico di lavoro nel periodo di stallo per privilegiare attività prioritarie di gestione e comunicazione, a causa degli impegni universitari del team.} \\
	\textbf{VI-20.2} & {Assegnazione di attività di 'palestra' per l'Analista, finalizzate allo studio delle correzioni AdR e alla stesura di quesiti per la proponente.} \\
	\textbf{VI-20.3} & {Divisione del ruolo\textsubscript{G} di Amministratore\textsubscript{G} in tre flussi di lavoro separati per ottimizzare la stesura dei documenti (PdP, PdQ), le configurazioni di CI/CD e le comunicazioni esterne.}
\end{longtable}

% --- 5. AZIONI (TODO) ---
\section{Azioni da Svolgere (To-Do)}

\renewcommand{\arraystretch}{1.3}
\begin{longtable}{p{2.3cm} c p{4.5cm} p{3cm} l}
    \toprule
    \textbf{Codice} & \textbf{Rif. Decisione} & \textbf{Descrizione Task\textsubscript{G}} & \textbf{Assegnatario} & \textbf{Scadenza} \\
    \midrule
    \endfirsthead

    \toprule
    \textbf{Codice} & \textbf{Rif. Decisione} & \textbf{Descrizione Task} & \textbf{Assegnatario} & \textbf{Scadenza} \\
    \midrule
    \endhead

    \midrule
    \multicolumn{5}{r}{\textit{Continua nella pagina successiva}} \\
    \midrule
    \endfoot

    \bottomrule
    \endlastfoot

    \ghissue{271} & - & Stesura del Diario di Bordo e del Verbale Interno odierno (VI N. 20) & Armando M. Scarati & 2026-07-07 \\[4pt]
	
	\ghissue{273} & VI-17.3 & Aggiornamento del PdQ e stesura e-mail di allineamento per Zucchetti & Roberto M. Doroftei & 2026-07-07 \\[4pt]

    \ghissue{276} & VI-17.3 & Stesura e-mail di allineamento per Zucchetti & Roberto M. Doroftei & 2026-07-07 \\[4pt]
	
	\ghissue{272} & VI-17.2 & Palestra (non rendicontata): studio correzioni AdR e quesiti proponente & Sofia De Blasi & 2026-07-07 \\[4pt]
	
	\ghissue{275} & VI-17.3 & Configurazione\textsubscript{G} iniziale CI/CD e definizione delle metriche di codice/test & Filippo Idiometri & 2026-07-07 \\[4pt]
	
	\ghissue{274} & VI-17.3 & Aggiornamento del Piano di Progetto\textsubscript{G} (PdP) & Francesco Marcon & 2026-07-07 \\[4pt]
\end{longtable}

% --- 6. PROSSIMA RIUNIONE ---
\section{Prossima Riunione}
\begin{itemize}
	\item \textbf{Data:} 2026-07-07
	\item \textbf{Ora:} 14:30
	\item \textbf{OdG preliminare\textsubscript{G}:}
	\begin{itemize}
		\item Verifica dei task assegnati per lo sprint corrente.
		\item Discussione sull'esito della comunicazione con Zucchetti.
		\item Aggiornamenti sullo stato di automazioni CI/CD e metriche di test.
		\item Pianificazione dello sprint successivo e assegnazione dei ruoli.
	\end{itemize}
\end{itemize}

